\origpage{201--227}

\BellaSection{7. Propriétés des bases}

Afin de s'assurer que la notion de base sert à quelque chose, il est bon d'établir quelques propriétés utiles. La première c'est la décomposition unique dans une base. On va également voir que les bases servent à caractériser les applications linéaires.

\medskip\noindent\textsc{Théorème 6.75.} --- \emph{Soit $\mathcal B=(e_i)_{i\in I}$ une base d'un espace vectoriel $E$ sur le corps $K$. Pour chaque vecteur $X$ de $E$ il existe une et une seule famille $(x_i)_{i\in I}$ de scalaires de $K$, presque tous nuls, telle que
\[
X=\sum_{i\in I}x_i e_i.
\]}

\medskip\noindent\textbf{6.75.bis} Rappelons que famille $(x_i)_{i\in I}$ de scalaires presque tous nuls, signifie que le cardinal de l'ensemble des $i$ tel que $x_i\ne0$ est fini. C'est forcément le cas en dimension finie, sur de tels espaces le presque tous nuls est inutile.

Il faut par ailleurs remarquer que la notation $\sum_{i\in I}x_i e_i$ représente une somme finie, que l'on pourrait indexer par les indices $i$ en nombre fini tels que $x_i\ne0$. L'ennui c'est qu'il faudrait changer d'ensemble d'indices en changeant de vecteurs, d'où l'utilisation de ce presque tous nuls.

La justification du théorème est immédiate.

Soit $X$ de $E=\operatorname{Vect}(\mathcal B)$, donc (Théorème 6.19), il existe une partie finie de $\mathcal B$, $\{e_{i_1},e_{i_2},\ldots,e_{i_n}\}$ par exemple, telle que $X$ soit combinaison linéaire de $e_{i_1},\ldots,e_{i_n}$ : d'où des scalaires $x_{i_1},\ldots,x_{i_n}$ tels que
\[
X=\sum_{k=1}^n x_{i_k}e_{i_k}.
\]
En posant $x_i=0$, $\forall i\ne i_1,\ne i_2,\ldots,\ne i_n$, on définit bien une famille $(x_i)_{i\in I}$ de scalaires de $K$ presque tous nuls telle que
\[
X=\sum_{i\in I}x_i e_i.
\]

Si on avait 2 écritures possibles,
\[
X=\sum_{i\in I}x_i e_i=\sum_{i\in I}x'_i e_i,
\]
on en déduit que
\[
0=\sum_{i\in I}(x_i-x'_i)e_i,
\]
les sommes portant \emph{a priori} sur un ensemble fini d'indices, ceux associés à des $x_i\ne0$ et à des $x'_i\ne0$ : pour préciser les choses on peut noter
\[
I_1=\{i,x_i\ne0\},\qquad I_2=\{i,x'_i\ne0\},\qquad I_3=I_1\cup I_2
\]
est fini, et on ne garde que
\[
X=\sum_{i\in I_3}x_i e_i=\sum_{i\in I_3}x'_i e_i
\quad\Longrightarrow\quad
\sum_{i\in I_3}(x_i-x'_i)e_i=0.
\]
D'où, $I_3$ étant une partie finie de $I$ et $\mathcal B$ une famille libre, on a $x_i-x'_i=0$, $\forall i\in I_3$. Comme pour $i\notin I_3$, $x_i=0=x'_i$ on a bien $x_i=x'_i$, $\forall i\in I$. \bookend

\medskip\noindent\textsc{Définition 6.76.} --- \emph{Les scalaires $x_i$ ainsi introduits sont les coordonnées de $X$ dans la base $\mathcal B$.}

Un vecteur est donc caractérisé par la donnée de ses coordonnées dans une base.

Passons aux applications linéaires.

On va voir que si $\mathcal B=(e_i)_{i\in I}$ est une base de $E$, la donnée de $u$ application linéaire de $E$ dans $F$ équivaut à la donnée des images $u(e_i)$ des vecteurs de cette base.

\medskip\noindent\textsc{Théorème 6.77.} --- \emph{Soient $E$ et $F$ deux espaces vectoriels sur un corps $K$, $\mathcal B=(e_i)_{i\in I}$ une base de $E$ et $(Y_i)_{i\in I}$ une famille de vecteurs de $F$ indexée par le même ensemble d'indices. Alors il existe une et une seule application linéaire $u$ de $E$ dans $F$ telle que, $\forall i\in I$, $u(e_i)=Y_i$. De plus on a
\[
\begin{aligned}
u\text{ injective}&\Longleftrightarrow (Y_i)_{i\in I}\text{ est libre},\\
u\text{ surjective}&\Longleftrightarrow (Y_i)_{i\in I}\text{ est génératrice de }F,\\
u\text{ isomorphisme}&\Longleftrightarrow (Y_i)_{i\in I}\text{ est une base de }F.
\end{aligned}
\]}

Existence et unicité de $u$ : il faut pouvoir définir $u(X)$ pour un $X$ quelconque de $E$.

Comme $\mathcal B$ est une base, à $X$ on associe la famille unique de ses coordonnées dans la base et si
\[
X=\sum_{i\in I}x_i e_i
\]
avec des $x_i$ presque tous nuls comme $u$ doit être linéaire, on doit avoir forcément
\[
u(X)=\sum_{i\in I}x_i u(e_i)=\sum_{i\in I}x_iY_i
\]
vu les conditions imposées sur $u$.

Il y a donc une et une seule définition possible de $u$. Le problème est de vérifier que $u$ ainsi définie est linéaire, or si on a $X'=\sum_{i\in I}x'_i e_i$, et si $\lambda$ et $\lambda'$ sont 2 scalaires, on aura
\[
\lambda X+\lambda'X'=\sum_{i\in I}(\lambda x_i+\lambda' x'_i)e_i,
\]
avec en fait un nombre fini seulement de $\lambda x_i+\lambda' x'_i$ non nuls; on posera donc
\[
u(\lambda X+\lambda'X')=\sum_{i\in I}(\lambda x_i+\lambda' x'_i)Y_i,
\]
expression qui s'écrit en fait
\[
\lambda\sum_{i\in I}x_iY_i+\lambda'\sum_{i\in I}x'_iY_i
\]
soit $\lambda u(X)+\lambda'u(X')$.

On a bien justifié l'existence et l'unicité de $u$ linéaire.

On a alors si $u$ injective, si $\{i_1,\ldots,i_n\}$ est une partie finie de $I$ et si $\lambda_1,\ldots,\lambda_n$ sont des scalaires tels que
\[
\sum_{k=1}^n\lambda_kY_{i_k}=0,
\]
c'est encore
\[
u\left(\sum_{k=1}^n\lambda_ke_{i_k}\right)=0,
\]
donc
\[
X=\sum_{k=1}^n\lambda_ke_{i_k}\in\operatorname{Ker} u=\{0\}
\]
(injectivité de $u$) donc $X$ est nul, or les $\lambda_k$ apparaissent comme ses coordonnées dans la base $\mathcal B$, donc les $\lambda_k$ sont nuls : on a justifié que la seule combinaison linéaire nulle des $(Y_{i_k})_{1\le k\le n}$ est celle avec des coefficients nuls : la famille $\{Y_{i_1},\ldots,Y_{i_n}\}$ est libre, ceci est vrai pour toute partie extraite de $I$ : finalement $(Y_i)_{i\in I}$ est libre.

Réciproquement on suppose la famille des $(Y_i)_{i\in I}$ libre.

Si $X=\sum_{i\in I}x_i e_i$ est dans $\operatorname{Ker} u$, en notant $\{i_1,\ldots,i_n\}$ la partie finie de $I$ associée, \emph{a priori}, à des $x_i$ non nuls, on aura
\[
u(X)=0=\sum_{k=1}^n x_{i_k}Y_{i_k},
\]
or $(Y_i)_{i\in I}$ est libre, donc les $x_{i_k}$ sont tous nuls, et finalement $X$ est nul : on a $u$ injective.

Donc $u$ injective $\Longleftrightarrow (Y_i)_{i\in I}$ libre.

On a $u$ surjective $\Longleftrightarrow (Y_i)_{i\in I}$ génératrice, en effet
\[
(Y_i)_{i\in I}\text{ génératrice}\Longleftrightarrow \operatorname{Vect}(u(e_i)_{i\in I})=F
\]
or l'élément générique de $\operatorname{Vect}(u(e_i)_{i\in I})$ est $\sum_{i\in I}x_i u(e_i)$, les $x_i$ étant presque tous nuls, soit encore $u(\sum_{i\in I}x_i e_i)$ par linéarité de $u$. Comme $\sum_{i\in I}x_i e_i$ est l'élément générique de $E$, on a donc $\operatorname{Vect}(u(e_i)_{i\in I})=u(E)$ en fait, d'où l'équivalence.

En réunissant ces 2 caractérisations, on a celle des bijections.

Le théorème va nous permettre de mieux connaître l'espace vectoriel $\mathcal L(E,F)$, si on connaît des bases de $E$ et de $F$.

\medskip\noindent\textsc{Théorème 6.78.} --- \emph{Soient $E$ et $F$ deux espaces vectoriels sur un corps $K$, $\mathcal B=(e_j)_{j\in J}$ une base de $E$ et $\mathcal C=(\varepsilon_i)_{i\in I}$ une base de $F$. Les applications $(u_{i,j})_{(i,j)\in I\times J}$ linéaires, définies pour chaque couple $(i,j)$ de $I\times J$ par les conditions :
\[
u_{i,j}(e_j)=\varepsilon_i\qquad\text{et}\qquad \forall j'\ne j,\ j'\in J,\quad u_{i,j}(e_{j'})=0
\]
forment une famille libre de $\mathcal L(E,F)$. C'est une base de $\mathcal L(E,F)$ si et seulement si $E$, (espace de départ) est de dimension finie.}

Il faut bien comprendre que pour un couple $(i,j)$ fixé de $I\times J$, $u_{i,j}$ est la seule application linéaire qui envoie tous les vecteurs de la base $\mathcal B$ de $E$ sur 0 sauf le $e_j$ qui est envoyé sur $\varepsilon_i$.

La façon d'indexer les applications servira lors du calcul matriciel. Le théorème 6.77 justifie l'existence des $u_{i,j}$ linéaires de $E$ dans $F$ vérifiant les conditions données.

Cette famille est libre, car si on prend un nombre fini des couples $(i,j)$, disons
\[
(i_1,j_1),(i_2,j_2),\ldots,(i_n,j_n),
\]
et si on a des scalaires $\lambda_1,\ldots,\lambda_n$ tels que l'application linéaire
\[
\sum_{k=1}^n\lambda_k u_{i_k,j_k}
\]
soit l'élément nul de $\mathcal L(E,F)$, c'est-à-dire l'application linéaire identiquement nulle, et si $j$ est l'un des indices de l'ensemble $\{j_1,\ldots,j_n\}$, en particulier l'image de $e_j$ est nulle.

Attention, les $n$ couples $(i_k,j_k)$ sont distincts, mais pas forcément les $n$ indices $j_1,\ldots,j_n$. Seulement, si $j\in\{j_1,\ldots,j_n\}$ et si pour plusieurs valeurs de $k$ on a $j_k=j$, les $i_k$ associés seront distincts. On peut donc écrire l'image de $e_j$ sous la forme :
\[
\left(\sum_{k=1}^n\lambda_k u_{i_k,j_k}\right)(e_j)
=
\sum_{\substack{1\le k\le n\\j_k=j}}\lambda_k\varepsilon_{i_k}=0,
\]
(car si $j_k\ne j$, $u_{i_k,j_k}$ envoie $e_j$ sur 0 et si $j_k=j$, $u_{i_k,j}$ envoie $e_j$ sur $\varepsilon_{i_k}$), et comme les $i_k$, en nombre fini, qui restent sont distincts, et que $\mathcal C=(\varepsilon_i)_{i\in I}$ est une base de $F$, les $\lambda_k$ associés aux $k$ tels que $j_k=j$ sont nuls. En faisant cela pour tout $j$ de $\{j_1,j_2,\ldots,j_n\}$ on obtient finalement la nullité de tous les $\lambda_k$, $k$ variant de 1 à $n$, et le fait que toute famille finie extraite de $(u_{i,j})_{(i,j)\in I\times J}$ est libre ce qui est la justification du fait que cette famille est libre.

Supposons alors $E$ de dimension finie, $n$, et soit $\mathcal B=\{e_1,\ldots,e_n\}$ une base de $E$, on va prouver que la famille des $(u_{i,j})$ est alors une base de $\mathcal L(E,F)$, c'est-à-dire qu'elle est génératrice. Soit donc $u\in\mathcal L(E,F)$.

Pour chaque $j$ compris entre 1 et $n$, $u(e_j)$ se décompose dans la base $\mathcal C=(\varepsilon_i)_{i\in I}$ de $F$ : il existe donc des scalaires presque tous nuls, que l'on note $x_{i,j}$, (l'indice $j$ pour rappeler qu'on travaille pour $e_j$), tels que
\[
u(e_j)=\sum_{i\in I}x_{i,j}\varepsilon_i.
\]
En fait, on peut introduire, pour chaque $j$ de $\{1,\ldots,n\}$, la partie finie $I_j$ de $I$ telle que
\[
x_{i,j}\ne0\Longleftrightarrow i\in I_j,
\]
(éventuellement $I_j$ est vide si $u(e_j)=0$).

Soit
\[
L=\bigcup_{j=1}^n I_j,
\]
c'est une partie finie de $I$, et si pour un $j$ fixé, on a $i\in L$ et $i\notin I_j$ on pose $x_{i,j}=0$. Il en résulte que
\[
\sum_{i\in L}x_{i,j}\varepsilon_i
=
\sum_{i\in I_j}x_{i,j}\varepsilon_i
=u(e_j),
\qquad(i\notin I_j\Rightarrow x_{i,j}=0)
\]
mais on a aussi $\varepsilon_i=u_{i,j}(e_j)$ d'où l'égalité
\[
u(e_j)=\sum_{i\in L}x_{i,j}u_{i,j}(e_j).
\]
Enfin, pour $j'\ne j$, chaque $u_{i,j'}(e_j)=0$, donc on a aussi
\[
\left[\sum_{j'=1}^n\left(\sum_{i\in L}x_{i,j'}u_{i,j'}\right)\right](e_j)
=
\sum_{i\in L}x_{i,j}u_{i,j}(e_j)
=
\sum_{i\in I_j}x_{i,j}\varepsilon_i
=u(e_j).
\]
La combinaison linéaire finie $\sum_{j'=1}^n\sum_{i\in L}x_{i,j'}u_{i,j'}$, d'une part et $u$ d'autre part donnent la même image de chaque vecteur $e_j$ de la base $\mathcal B$ de $E$ : elles coïncident d'après le théorème 6.77. On vient donc de vérifier que, si $\dim(E)$ finie la famille des $(u_{i,j})_{(i,j)\in I\times J}$ est une base de $\mathcal L(E,F)$ car libre et génératrice. \bookend

\emph{Remarque :} Si $\dim E$ infinie, cette famille libre des $(u_{i,j})$ n'est jamais une base, donc jamais génératrice, car chaque $u_{i,j}$ envoyant toute la base $\mathcal B$ sur 0 sauf $e_j$ sur $\varepsilon_i$ une combinaison linéaire d'un nombre fini des $u_{i,j}$ envoie tous les vecteurs de $\mathcal B$ sur 0 sauf un nombre fini. Il en résulte que par exemple l'application $u$ envoyant chaque $e_j$, (en nombre infini) sur un $\varepsilon_{i_0}$ fixé, ($u$ existe d'après le Théorème 6.77, appliqué à la base $\mathcal B=(e_j)_{j\in J}$ de $E$ et la famille constante $(Y_j)_{j\in J}$ de $F$ avec $\forall j\in J$, $Y_j=\varepsilon_{i_0}$), ne peut pas être dans $\operatorname{Vect}((u_{i,j})_{(i,j)\in I\times J})$.

\emph{Cas particulier de $E$ et $F$ tous les deux de dimension finie.}

\medskip\noindent\textsc{Corollaire 6.79.} --- \emph{Si $E$ est de dimension finie, $p$, et $F$ de dimension finie, $n$, l'espace vectoriel $\mathcal L(E,F)$ est de dimension finie $np$.}

Car avec $\mathcal B=(e_j)_{1\le j\le p}$ base de $E$ et $\mathcal C=(\varepsilon_i)_{1\le i\le n}$ base de $F$, il y a alors $np$ applications du type $u_{i,j}$, qui forment une base de $\mathcal L(E,F)$ qui est donc de dimension finie $np$. \bookend

Les deux théorèmes d'isomorphisme, (Théorèmes 6.41 et 6.43) vont nous fournir des égalités intéressantes, quand on appliquera le résultat suivant, valable en fait en dimension quelconque :

\medskip\noindent\textsc{Théorème 6.80.} --- \emph{Deux espaces vectoriels $E$ et $F$ sont isomorphes si et seulement si ils ont des bases équipotentes.}

On dira que les espaces sont isomorphes s'il y a un isomorphisme de l'un sur l'autre.

En effet si $E$ et $F$ sont isomorphes, soit $u$ un isomorphisme de $E$ sur $F$, $\mathcal B=(e_j)_{j\in J}$ une base de $E$, d'après le théorème 6.77, $u$ est la seule application linéaire envoyant chaque $e_j$ sur $u(e_j)$ \BellaCorrectionNote{vecteur de $F$}{vecteur de $E$}{Puisque $u:E\to F$, l'image $u(e_j)$ appartient à $F$.} et $u$ étant bijective, cette famille $\mathcal C=(u(e_j))_{j\in J}$ est une base de $F$ qui admet donc des bases équipotentes à celles de $E$.

\emph{Réciproque}

Si $E$ et $F$ admettent des bases équipotentes, on peut les indexer par le même ensemble d'indices, $J$, donc soit $\mathcal B=(e_j)_{j\in J}$ une base de $E$ et $\mathcal C=(\varepsilon_j)_{j\in J}$ une base de $F$, toujours d'après le théorème 6.77, il existe une et une seule application linéaire envoyant chaque $e_j$ sur $\varepsilon_j$, cette application étant bijective puisque $\mathcal C$ est une base de $F$. On a bien $E$ et $F$ isomorphes. \bookend

De même, quelque soit la dimension des espaces on peut établir un lien entre somme directe et bases.

\medskip\noindent\textsc{Théorème 6.81.} --- \emph{Soit un espace vectoriel $E$ somme directe de $n$ sous-espaces $F_1,F_2,\ldots,F_n$ et soit, pour chaque $i=1,2,\ldots,n$, $\mathcal B_i$ une base de $F_i$, alors $\mathcal B=\bigcup_{i=1}^n\mathcal B_i$ est une base de $E$, les $\mathcal B_i$ étant 2 à 2 disjointes.}

Car si $X\in E$, $X$ s'écrit dans la somme directe $E=\bigoplus_{i=1}^n F_i$ sous la forme $X=X_1+\cdots+X_n$ avec chaque $X_i$ dans $F_i=\operatorname{Vect}\mathcal B_i$, donc $X_i$ est combinaison linéaire d'un nombre fini d'éléments de $\mathcal B_i$, donc de $\mathcal B$. En sommant pour $i$ variant de 1 à $n$, on obtient $X$ combinaison linéaire d'un nombre fini d'éléments de $\mathcal B$.

Il est clair que si $i\ne j$, $\mathcal B_i\cap\mathcal B_j=\varnothing$ sinon, si $e\in\mathcal B_i\cap\mathcal B_j$ on aurait $e\ne0$ dans
\[
F_i\cap\left(\sum_{\substack{k=1\\k\ne i}}^n F_k\right),
\]
(car $e$ dans $F_j\subset\sum_{k\ne i}F_k$) ce qui est exclu.

$\mathcal B$ est enfin partie libre de $E$, car si on a une combinaison d'un nombre fini des éléments de $\mathcal B$, nulle, soit $X$ cette combinaison linéaire. On peut noter
\[
\mathcal B_1=\{e_{i_1},i_1\in I_1\},\quad \mathcal B_2=\{e_{i_2},i_2\in I_2\},\quad\ldots,\quad \mathcal B_n=\{e_{i_n},i_n\in I_n\}
\]
les bases de $F_k$, les ensembles d'indices $I_1,\ldots,I_n$ étant disjoints, puis introduire les parties finies suivantes :
\[
J_1=\{i_1\in I_1,e_{i_1}\text{ intervient avec le coefficient }x_{i_1}\text{ dans }X\},
\]
\[
J_2=\{i_2\in I_2,e_{i_2}\text{ intervient avec le coefficient }x_{i_2}\text{ dans }X\},
\]
et enfin
\[
J_n=\{i_n\in I_n,e_{i_n}\text{ intervient avec le coefficient }x_{i_n}\text{ dans }X\};
\]
alors
\[
X=\sum_{k=1}^n\left(\sum_{i_k\in J_k}x_{i_k}e_{i_k}\right)=0.
\]
Si on pose $X_k=\sum_{i_k\in J_k}x_{i_k}e_{i_k}$, l'écriture $X_1+\cdots+X_n=0$ avec les $X_j\in F_j$ et $E=\bigoplus_{j=1}^nF_j$ implique que chaque $X_k=0$, (théorème 6.51), mais alors chaque $x_{i_k}$ est nul, $X_k$ étant décomposé dans une base $\mathcal B_k$ de $F_k$.

Finalement $\mathcal B$ est bien libre et génératrice : c'est une base de $E$. \bookend

\medskip\noindent\textsc{Théorème 6.82.} --- \emph{Soit $\mathcal B=(e_i)_{i\in I}$ une base d'un espace vectoriel $E$ et $I=\bigcup_{k=1}^n I_k$ une partition de $I$. En posant $\mathcal B_k=(e_i)_{i\in I_k}$ et $F_k=\operatorname{Vect}\mathcal B_k$, on a
\[
E=\bigoplus_{k=1}^n F_k.
\]}

Car tout $X$ de $E$ se décompose en combinaison d'un nombre fini des $e_i$.

En reprenant les notations de la justification précédente, avec $J_1$ ensemble des $i_1$ de $I_1$ tels que $e_{i_1}$ intervient avec le coefficient $x_{i_1}$ dans $X$, ... on arrive à
\[
X=\sum_{k=1}^n X_k\qquad\text{avec}\qquad X_k=\sum_{i_k\in J_k}x_{i_k}e_{i_k}\in F_k.
\]
Donc on a déjà composé chaque $X$ de $E$ sous la forme $\sum_{k=1}^nX_k$ avec les $X_k\in F_k$.

Puis si une telle somme est nulle, en décomposant $X_k$ dans la base $\mathcal B_k$ de $F_k$, sous la forme
\[
X_k=\sum_{i_k\in I_k}x_{i_k}e_{i_k},
\]
les $x_{i_k}$ étant presque tous nuls, on a
\[
X=\sum_{k=1}^nX_k=0
=\sum_{k=1}^n\left(\sum_{i_k\in I_k}x_{i_k}e_{i_k}\right).
\]
C'est une combinaison linéaire nulle d'un nombre fini des $e_i$ de $\mathcal B=\bigcup_{k=1}^n\mathcal B_k$, base donc famille libre : tous les $x_{i_k}$ sont nuls d'où chaque $X_k=0$. On a bien vérifié l'une des conditions de somme directe (Théorème 6.51). \bookend

On applique tout ce qui précède aux espaces vectoriels de dimension finie. On a

\medskip\noindent\textsc{Corollaire 6.83.} --- \emph{Soient $G$ et $F$ deux sous-espaces supplémentaires d'un espace vectoriel $E$ de dimension finie : on a
\[
\dim E=\dim F+\dim G.
\]}

Car une base de $E$ est obtenue en réunissant une base de $F$ et une de $G$, ces bases étant disjointes : c'est le théorème 6.81. \bookend

\medskip\noindent\textsc{Corollaire 6.84.} --- \emph{Soit $F$ un sous-espace vectoriel de $E$ espace de dimension finie : on a
\[
\dim(E/F)=\dim E-\dim F.
\]}

Car $F$ admet des supplémentaires (Théorème 6.69). Soit $G$ un supplémentaire de $F$ et $p$ la projection sur $G$ de noyau $F$. Les espaces vectoriels $E/\operatorname{Ker} p$ et $p(E)$ sont isomorphes (premier théorème d'isomorphisme, théorème 6.41), donc de même dimension. Comme $\operatorname{Ker} p=F$ et $p(E)=G$ on a finalement
\[
\dim(E/F)=\dim G=\dim E-\dim F
\]
vu le corollaire 6.83 puisque $E=F\oplus G$. \bookend

\medskip\noindent\textsc{Corollaire 6.85.} --- \emph{Soit $E$ un espace vectoriel de dimension finie, $F$ un espace vectoriel de dimension quelconque, et $f$ linéaire de $E$ dans $F$. On a
\[
\dim(E)=\dim(f(E))+\dim(\operatorname{Ker} f).
\]}

Car les espaces vectoriels $E/\operatorname{Ker} f$ et $f(E)$ sont isomorphes, (premier théorème d'isomorphisme, théorème 6.41) donc de même dimension (Théorème 6.80) et on applique le corollaire 6.84 pour calculer $\dim(E/\operatorname{Ker} f)$ d'où
\[
\dim E-\dim(\operatorname{Ker} f)=\dim(f(E))
\]
ce qui est le résultat. \bookend

\textsc{Conséquence} si $E$ et $F$ sont de même dimension finie, on a, pour $f\in\mathcal L(E,F)$ :
\[
(f\text{ injective})\Longleftrightarrow(f\text{ surjective})\Longleftrightarrow(f\text{ bijective}).
\]

\medskip\noindent\textsc{Définition 6.86.} --- \emph{Soient $E$ et $F$ deux espaces vectoriels, $f\in\mathcal L(E,F)$. Lorsque $f(E)$ est un sous-espace de dimension finie de $F$, sa dimension est appelée rang de $f$.}

\medskip\noindent\textsc{Corollaire 6.87.} --- \emph{Soient $F$ et $G$ deux sous-espaces d'un espace vectoriel $E$ de dimension finie : on a
\[
\dim(F+G)=\dim F+\dim G-\dim(F\cap G).
\]}

C'est la traduction, en dimension, du 2e théorème d'isomorphisme, (Théorème 6.43), qui nous dit, (si tant est qu'il ait la parole), que
\[
(F+G)/F\simeq G/(F\cap G)
\]
d'où en prenant les dimensions (6.80 et 6.84) :
\[
\dim(F+G)-\dim F=\dim G-\dim(F\cap G)
\]
ce qui est la formule voulue. \bookend

Ces quelques corollaires sont des outils bien commodes pour l'étude des situations rencontrées dans les espaces vectoriels. Avant d'aborder un paragraphe bien plus mystérieux, il reste une remarque à faire :

\medskip\noindent\textsc{Remarque 6.88.} --- Soit $E$ un espace vectoriel de dimension finie.
\begin{enumerate}[label=\arabic*),leftmargin=2.1em]
\item Si on a $F$ sous-espace de $E$ et $\dim F=\dim E$ alors $F=E$.
\item Si $G\subset E$ avec $G\simeq E$ alors $G=E$.
\end{enumerate}
Mais ces 2 résultats sont faux en dimension infinie où on peut rencontrer $F$ sous-espace strict de $E$, $F$ et $E$ étant de même dimension, (donc isomorphes).

En dimension finie, ($n=\dim E$) si $\mathcal B$ est une base de $F$ et si $\operatorname{card}\mathcal B=n$, (c'est le cas si $\dim F=\dim E$ ou si $F$ est isomorphe à $E$) le théorème de la base incomplète (Théorème 6.66) nous dit que $\mathcal B$ est base de $E$ donc $F$ (ou $G$) $=E$.

En dimension infinie, reprenons $E=K[X]$ espace des polynômes sur $K$. On note
\[
e_0=(1,0,0,\ldots,0,\ldots)
\]
la suite où tous les termes sont nuls sauf le premier et plus généralement
\[
e_p=(0,\ldots,0,1,0,\ldots),\qquad(1\text{ au rang }p+1).
\]
On peut vérifier que $(e_n)_{n\in\mathbb N}$ est une base de $E$.

On a $F=\operatorname{Vect}(e_{2n})_{n\in\mathbb N}$ est sous-espace strict de $E$, (on a les polynômes pairs) et pourtant les bases $(e_n)_{n\in\mathbb N}$ et $(e_{2n})_{n\in\mathbb N}$ sont équipotentes d'où $F$ isomorphe à $E$ avec $F\subsetneq E$.


\BellaSection{8. Dualité, espace dual}

\medskip\noindent\textsc{Définition 6.89.} --- \emph{Soit $E$ un espace vectoriel sur un corps $K$, on appelle espace vectoriel dual de $E$, l'espace vectoriel $\mathcal L(E,K)$ des applications linéaires de $E$ dans le corps $K$.}

On note encore $E^*$ cet espace vectoriel et on appelle forme linéaire tout élément de $E^*$, c'est-à-dire toute application linéaire de $E$ dans le corps $K$.

\medskip\noindent\textbf{6.90.} Si $f$ est un élément non nul du dual, son noyau est appelé hyperplan de $E$. C'est un sous-espace de $E$ dont les supplémentaires sont de dimension 1.

En effet, soit $H=\operatorname{Ker} f$ le noyau de $f$ élément non nul du dual. Le premier théorème d'isomorphisme (Théorème 6.41) nous dit que $E/\operatorname{Ker} f$ est isomorphe à $f(E)$. Mais $f(E)$ est sous-espace vectoriel du corps $K$, non réduit à $\{0\}$ car $f\ne0$, donc $f(E)=K$, (si on a $\lambda\ne0$ dans $f(E)$, $\forall\mu\in K$, $\mu=(\lambda^{-1}\mu)\lambda$ donc si le vecteur $X$ de $E$ est tel que $f(X)=\lambda$ on a $f(\lambda^{-1}\mu\cdot X)=\mu$). Or $K$ est espace vectoriel de dimension 1 sur lui-même, le scalaire $1_K$ (élément neutre de $K$ pour le produit) formant une base.

Soit alors $L$ un supplémentaire de $H$ dans $E$, (il y en a, théorème 6.73), en utilisant la projection sur $L$, de noyau $H$, on sait aussi que $E/H=E/\operatorname{Ker} f\simeq L$, d'où $L$ et $f(E)$ isomorphes donc de bases équipotentes : on a bien $L$ de dimension 1. En fait on a la réciproque : si un sous-espace $H$ de $E$ admet pour supplémentaires des sous-espaces de dimension 1, c'est le noyau d'une forme linéaire.

Car soit $L$ un supplémentaire de $H$, $\mathcal B_1$ une base de $H$, $\{a\}$ une base de $L$, (de cardinal 1 par hypothèse), comme $E$ est somme directe de $L$ et $H$ on a $\mathcal B=\mathcal B_1\cup\{a\}$ base de $E$, (Théorème 6.81).

On définit alors $f$ linéaire de $E$ dans $K$ en se donnant l'image des vecteurs de la base $\mathcal B$, (Théorème 6.77) et pour cela on pose $f(a)=1$ et tout vecteur de $\mathcal B_1$ est envoyé sur 0.

Soit $X$ de $E$, décomposé en $X=Y+\lambda\cdot a$ avec $Y\in H$ et $\lambda$ scalaire, décomposition unique dans la somme directe $E=H\oplus L$. On a alors
\[
f(X)=f(Y)+\lambda f(a)=0+\lambda=\lambda
\]
donc
\[
X\in\operatorname{Ker} f\Longleftrightarrow\lambda=0\Longleftrightarrow X\in H.
\]
On a bien construit $f$ de noyau $H$. Il y a non unicité de $f$, d'abord parce que $\varphi$ envoyant $\mathcal B_1$ sur 0 et $a$ sur un scalaire $\ne0$ convient aussi, ensuite parce que $L$ n'est pas unique. \bookend

\medskip\noindent\textsc{Définition 6.91.} --- \emph{On appelle codimension d'un sous-espace vectoriel $F$ de l'espace $E$ la dimension de $E/F$, ou, ce qui est équivalent, la dimension des supplémentaires de $F$ dans $E$.}

Bien sûr, dans cette définition, si $E/F$ est de dimension infinie, la codimension devient un cardinal infini.

Avec cette définition on a établi le

\medskip\noindent\textsc{Théorème 6.92.} --- \emph{Soit $E$ un espace vectoriel sur un corps $K$, un sous-espace $H$ est un hyperplan de $E$ si et seulement si il est de codimension 1.} \bookend

Au fait, pourquoi hyperplan? Parce que, historiquement, les espaces vectoriels ont été abordés par l'étude de la géométrie, et de l'espace de dimension 3, espace « usuel » dans lequel les sous-espaces de codimension 1 sont de dimension $3-1=2$, donc sont les plans. La généralisation de cette notion fait passer aux hyperplans, comme on passe des marchés aux hypermarchés! C'est du langage hyperbolique.

\medskip\noindent\textbf{Orthogonalité}

\medskip\noindent\textsc{Définition 6.93.} --- \emph{Une forme linéaire $f$, $f\in E^*$, et un vecteur $X$ de $E$ sont dits orthogonaux si et seulement si $f(X)=0$.}

Donc c'est, pour $f\ne0$, traduire que $X$ est dans l'hyperplan $\operatorname{Ker} f$.

\medskip\noindent\textsc{Définition 6.94.} --- \emph{Soit $A$ une partie de l'espace vectoriel $E$. On appelle orthogonal de $A$, l'ensemble des $f$ du dual tels que, $\forall X\in A$, $f(X)=0$.}

On le note $A^\perp$, on a donc
\[
A^\perp=\{f;\ f\in E^*,\ \forall X\in A,\ f(X)=0\}.
\]

En remarquant que, pour un hyperplan $H$ de $E$, tout $f$ du dual tel que $\operatorname{Ker} f=H$ s'appelle encore une équation de l'hyperplan, on peut interpréter $A^\perp$ comme l'ensemble des équations des hyperplans contenant la partie $A$, ensemble auquel on adjoint la forme nulle.

\medskip\noindent\textsc{Théorème 6.95.} --- \emph{Soit $A$ une partie de l'espace vectoriel $E$, alors $A^\perp$ est un sous-espace vectoriel de $E^*$. Si $A\subset B$, $B^\perp\subset A^\perp$. On a $A^\perp=(\operatorname{Vect}A)^\perp$.}

D'abord $A^\perp$ non vide car la forme identiquement nulle annule tout vecteur de $A$, donc est dans $A^\perp$, puis $A^\perp$ stable par combinaison linéaire : si $f$ et $g$ sont dans $A^\perp$, $\lambda$ et $\mu$ dans le corps, alors $\forall X\in A$ on a
\[
(\lambda f+\mu g)(X)=\lambda f(X)+\mu g(X)=0
\]
puisque $f(X)$ et $g(X)$ sont nuls.

Il est clair que si $A\subset B$, et si un hyperplan contient $B$, alors il contient $A$ donc les équations des hyperplans contenant $B$, sont parmi les équations des hyperplans contenant $A$, soit $B^\perp\subset A^\perp$.

Enfin, pour un hyperplan (sous-espace vectoriel) contenir $A$ équivaut à contenir $\operatorname{Vect}A$ d'où en passant aux équations $A^\perp=(\operatorname{Vect}A)^\perp$. \bookend

\medskip\noindent\textsc{Théorème 6.96.} --- \emph{Soient $F$ et $G$ deux sous-espaces vectoriels de $E$. On a
\[
(F+G)^\perp=F^\perp\cap G^\perp
\qquad\text{et}\qquad
(F\cap G)^\perp=F^\perp+G^\perp.
\]}

La première égalité est facile à justifier : $F\subset(F+G)$ et $G\subset(F+G)$ impliquent $(F+G)^\perp\subset F^\perp$ et $(F+G)^\perp\subset G^\perp$ d'où $(F+G)^\perp\subset(F^\perp\cap G^\perp)$; puis si $f\in F^\perp\cap G^\perp$, pour $X$ de $(F+G)$ écrit sous la forme $y+z$ avec $y$ dans $F$ et $z$ dans $G$ on a $f(X)=f(y)+f(z)$, (linéarité). Or $f$ est dans $F^\perp$ et $y$ dans $F$ d'où $f(y)=0$ et $f$ dans $G^\perp$, $z$ dans $G$ donne $f(z)=0$, donc $f(X)=0$ : on a bien $f\in(F+G)^\perp$; d'où l'égalité $(F+G)^\perp=F^\perp\cap G^\perp$.

Pour la deuxième égalité on a : $(F\cap G)\subset F$ et $(F\cap G)\subset G$ d'où $F^\perp$ et $G^\perp$ sous-espaces de $(F\cap G)^\perp$ donc \emph{a fortiori}
\[
F^\perp+G^\perp=\operatorname{Vect}(F^\perp\cup G^\perp)\subset(F\cap G)^\perp.
\]

Pour l'autre inclusion, soit $\varphi\in(F\cap G)^\perp$ que l'on veut décomposer en $\varphi_1+\varphi_2$ avec $\varphi_1\in F^\perp$ et $\varphi_2\in G^\perp$. Pour cela, on part de $\mathcal B_1$ base de $F\cap G$, il existe $\mathcal C_1$ tel que $\mathcal B_1\cup\mathcal C_1$ soit une base de $F$ et il existe $\mathcal C_2$ tel que $\mathcal B_1\cup\mathcal C_2$ soit une base de $G$, ceci par application adéquate du théorème de la base incomplète : théorème 6.72.

On a $\mathcal B_1\cup\mathcal C_1\cup\mathcal C_2$ famille libre car si on a des scalaires $\lambda_i$, $\mu_j$, $\rho_k$ presque tous nuls tels que
\[
\sum_{X_i\in\mathcal B_1}\lambda_iX_i+\sum_{Y_j\in\mathcal C_1}\mu_jY_j+\sum_{Z_k\in\mathcal C_2}\rho_kZ_k=0,
\]
en posant
\[
X=\sum_{X_i\in\mathcal B_1}\lambda_iX_i+\sum_{Y_j\in\mathcal C_1}\mu_jY_j
\qquad\text{et}\qquad
Z=\sum_{Z_k\in\mathcal C_2}\rho_kZ_k,
\]
on a $X$ dans $F$ et $Z$ dans $G$, de plus $Z=-X$ donc $Z\in F\cap G=\operatorname{Vect}(\mathcal B_1)$.

Or $\mathcal B_1\cup\mathcal C_2$ est une base de $G$, le vecteur $Z$ devant se décomposer uniquement en fonction des vecteurs de $\mathcal B_1$, c'est que ses coordonnées suivant les vecteurs de $\mathcal C_2$ sont nulles, d'où chaque $\rho_k=0$ mais alors $X=0$ avec $X$ décomposé dans la base $\mathcal B_1\cup\mathcal C_1$ de $F$ d'où chaque $\lambda_i=0$ et chaque $\mu_j=0$ : la famille $\mathcal B_1\cup\mathcal C_1\cup\mathcal C_2$ est bien libre : on la complète en $\mathcal B_1\cup\mathcal C_1\cup\mathcal C_2\cup\mathcal L$ base de $E$ et on définit
\[
\begin{array}{c|cccc}
&\mathcal B_1&\mathcal C_1&\mathcal C_2&\mathcal L\\\hline
\varphi_1&0&0&\varphi&\varphi\\
\varphi_2&0&\varphi&0&0
\end{array}
\]
au sens où les formes linéaires prennent sur chacun de ces ensembles de vecteurs les valeurs indiquées.

On a bien $\varphi_1$ nulle sur $\mathcal B_1\cup\mathcal C_1$ base de $F$ donc $\varphi_1\in F^\perp$, $\varphi_2$ nulle sur $\mathcal B_1\cup\mathcal C_2$ base de $G$ donc $\varphi_2\in G^\perp$, et $\varphi_1+\varphi_2=\varphi$ sur $\mathcal C_1$, sur $\mathcal C_2$, sur $\mathcal L$ et aussi sur $\mathcal B_1$ car alors $\varphi=0$.

On a donc décomposé $\varphi$ de $(F\cap G)^\perp$ en $\varphi_1+\varphi_2\in F^\perp+G^\perp$ d'où finalement l'égalité voulue. \bookend

\medskip\noindent\textbf{6.97. Interprétation géométrique :} les hyperplans contenant $F+G$ sont exactement ceux qui contiennent $F$ et $G$. Le deuxième point ne s'interprète pas puisque le $F^\perp+G^\perp$ ne correspond pas directement à des hyperplans.

\medskip\noindent\textbf{6.98.} De manière analogue, si $A$ est une partie du dual cette fois, on définit son orthogonal par
\[
A^\perp=\{X;\ X\in E,\ \forall f\in A,\ f(X)=0\}.
\]
C'est encore
\[
A^\perp=\bigcap_{f\in A}\operatorname{Ker} f
\]
ce qui permet de justifier facilement
\begin{enumerate}[label=\arabic*),leftmargin=2.1em]
\item que $A^\perp$ est sous-espace vectoriel de $E$, (intersection de sous-espaces);
\item que $A\subset B\Rightarrow B^\perp\subset A^\perp$, (dans $B^\perp$ on a déjà l'intersection des $\operatorname{Ker} f$ pour les $f$ de $A$, et on coupe encore, $B^\perp$ diminue donc ... ).
\item On peut vérifier que $(\operatorname{Vect}A)^\perp=A^\perp$, car $A\subset\operatorname{Vect}A$ donne déjà $(\operatorname{Vect}A)^\perp\subset A^\perp$, puis si $X\in A^\perp$ et $f\in\operatorname{Vect}A$, $f$ est combinaison linéaire d'un nombre fini d'éléments de $A$, soient $f_1,\ldots,f_n$ ces éléments, $\lambda_1,\ldots,\lambda_n$ les scalaires coefficients de la combinaison linéaire, on a alors
\[
f=\sum_{i=1}^n\lambda_i f_i\qquad\text{donc}\qquad f(X)=\sum_{i=1}^n\lambda_i f_i(X)=0,
\]
chaque $f_i(X)$ étant nul puisque $X$ est dans $A^\perp$ et les $f_i$ dans $A$. Le vecteur $X$ annule bien chaque $f$ de $\operatorname{Vect}A$ : $X\in(\operatorname{Vect}A)^\perp$ d'où l'inclusion $A^\perp\subset(\operatorname{Vect}A)^\perp$ et l'égalité.
\end{enumerate}

\medskip\noindent\textsc{Théorème 6.99.} --- \emph{Soit $S$ une partie de l'espace vectoriel $E$, on a
\[
(S^\perp)^\perp=\operatorname{Vect}S.
\]}

D'abord $S^\perp=(\operatorname{Vect}S)^\perp$ (Théorème 6.95), et si on prend un vecteur $X$ de $\operatorname{Vect}S$, pour tout $f$ de $(\operatorname{Vect}S)^\perp$ on a $f(X)=0$ donc $X$ est bien dans l'orthogonal de la partie $(\operatorname{Vect}S)^\perp$ du dual $E^*$, d'où l'inclusion
\[
\operatorname{Vect}S\subset((\operatorname{Vect}S)^\perp)^\perp.
\]
On va prouver l'autre inclusion en justifiant que, si $X\notin\operatorname{Vect}S$, alors $X\notin((\operatorname{Vect}S)^\perp)^\perp$.

Soit donc $\mathcal C$ une base de $\operatorname{Vect}(S)$, $X\notin\operatorname{Vect}(S)$ donc $\mathcal C\cup\{X\}$ est une famille libre (Théorème 6.61), car si $\mathcal C\cup\{X\}$ liée, on obtiendrait $X$ combinaison linéaire d'un nombre fini de vecteurs de $\mathcal C$.

On complète en $\mathcal B=\mathcal C\cup\{X\}\cup\mathcal L$ base de $E$, (Théorème 6.72, dit de la base incomplète) et on définit une forme linéaire $\mu$ en donnant les images des vecteurs de $\mathcal B$, (Théorème 6.77). Pour cela, $\mu$ envoie $\mathcal C$ sur 0, donc $\mu$ annule $\operatorname{Vect}(\mathcal C)=\operatorname{Vect}(S)$, soit $\mu\in(\operatorname{Vect}S)^\perp$; et on impose $\mu(X)=1$, et n'importe quoi pour $\mu(\mathcal L)$.

On a $\mu\in(\operatorname{Vect}S)^\perp$ et $\mu(X)=1$ donc $X\notin((\operatorname{Vect}S)^\perp)^\perp$.

On a bien justifié l'inclusion $((\operatorname{Vect}S)^\perp)^\perp\subset\operatorname{Vect}S$ et l'égalité souhaitée. \bookend

Une autre notion utile dans le dual est celle de transposition.

\medskip\noindent\textsc{Définition 6.100.} --- \emph{Soient $E$ et $F$ deux espaces vectoriels sur un corps $K$, $E^*$ et $F^*$ leurs duaux et $u\in\mathcal L(E,F)$. On appelle transposée de $u$ l'application notée ${}^tu$ de $F^*$ dans $E^*$ définie par :
\[
\forall\varphi\in F^*,\qquad {}^tu(\varphi)=\varphi\circ u,
\]
(${}^tu$ se lit transposée de $u$).}

Ceci a un sens car $u$ est linéaire de $E$ dans $F$ et $\varphi$ de $F$ dans le corps $K$ donc $\varphi\circ u$ est linéaire de $E$ dans $K$ : on a bien ${}^tu(\varphi)\in E^*$.

\medskip\noindent\textsc{Théorème 6.101.} --- \emph{L'application ${}^tu$ est linéaire, la transposition est linéaire et ${}^t(u\circ v)={}^tv\circ{}^tu$, (sous réserve que $u\circ v$ ait un sens).}

On a
\[
{}^tu(\lambda_1\varphi_1+\lambda_2\varphi_2)=(\lambda_1\varphi_1+\lambda_2\varphi_2)\circ u:
\]
c'est l'application de $E$ dans $K$ qui à un vecteur $X$ de $E$ associe $(\lambda_1\varphi_1+\lambda_2\varphi_2)(u(X))$ ce qui vaut, compte tenu de la structure vectorielle, $\lambda_1\varphi_1(u(X))+\lambda_2\varphi_2(u(X))$ soit
\[
[\lambda_1{}^tu(\varphi_1)+\lambda_2{}^tu(\varphi_2)](X).
\]
Ceci étant vrai $\forall X\in E$, c'est que
\[
{}^tu(\lambda_1\varphi_1+\lambda_2\varphi_2)=\lambda_1{}^tu(\varphi_1)+\lambda_2{}^tu(\varphi_2)
\]
d'où la linéarité de ${}^tu$.

L'application transposition : $u\longmapsto{}^tu$ est elle aussi linéaire de $\mathcal L(E,F)$ dans $\mathcal L(F^*,E^*)$.

Car si $u,v$ sont dans $\mathcal L(E,F)$ et $\lambda,\mu$ dans le corps $K$, ${}^t(\lambda u+\mu v)$ est définie par
\[
\varphi\longmapsto{}^t(\lambda u+\mu v)(\varphi)=\varphi\circ(\lambda u+\mu v)
\]
ce qui, compte tenu de la linéarité de $\varphi$, vaut $\lambda\varphi\circ u+\mu\varphi\circ v$ (calculer pour un $X$ quelconque de $E$ pour le vérifier).

C'est encore $\lambda{}^tu(\varphi)+\mu{}^tv(\varphi)$ soit $(\lambda{}^tu+\mu{}^tv)(\varphi)$, d'après cette fois la structure vectorielle de $\mathcal L(F^*,E^*)$.

L'égalité
\[
{}^t(\lambda u+\mu v)(\varphi)=(\lambda{}^tu+\mu{}^tv)(\varphi)
\]
étant valable pour tout $\varphi$ de $F^*$, c'est que ${}^t(\lambda u+\mu v)=\lambda{}^tu+\mu{}^tv$, ce qui est bien la linéarité de la transposition.

Enfin soient $E,F,G$ trois espaces vectoriels et $u$ linéaire de $E$ dans $F$, $v$ linéaire de $F$ dans $G$. Alors $v\circ u$ est linéaire de $E$ dans $G$ et, pour tout $\psi$ de $G^*$ on a
\[
\begin{aligned}
{}^t(v\circ u)(\psi)&=\psi\circ(v\circ u), &&\text{définition de }{}^t(v\circ u),\\
&=(\psi\circ v)\circ u, &&\text{par associativité du produit de composition},\\
&=({}^tv(\psi))\circ u,\\
&={}^tu({}^tv(\psi))=({}^tu\circ{}^tv)(\psi).
\end{aligned}
\]
C'est vrai pour tout $\psi$ de $G^*$ donc ${}^t(v\circ u)={}^tu\circ{}^tv$. \bookend

\medskip\noindent\textsc{Théorème 6.102.} --- \emph{Soient $E$ et $F$ deux espaces vectoriels et $u\in\mathcal L(E,F)$. On a
\[
\operatorname{Ker}({}^tu)=(u(E))^\perp.
\]}

Car $\varphi$ de $F^*$ est dans $\operatorname{Ker}({}^tu)$ si et seulement si ${}^tu(\varphi)=\varphi\circ u=0$, c'est-à-dire $\forall X\in E$, $\varphi(u(X))=0$ ce qui est la transcription de $\varphi$ orthogonal à tout vecteur de l'image $u(E)$ d'où l'équivalence. \bookend

\medskip\noindent\textsc{Corollaire 6.103.} --- \emph{On a $u$ surjective $\Longleftrightarrow {}^tu$ injective.}

En effet établissons d'abord un lemme.

\medskip\noindent\textsc{Lemme 6.104.} --- \emph{Soit $X$ un vecteur non nul de $E$, il existe des formes linéaires non nulles sur $X$.}

Car la famille libre $\{X\}$ se complète en $\{X\}\cup\mathcal C$ base $\mathcal B$ de $E$ et la forme coordonnée $\mu$ qui envoie $X$ sur 1 et tout vecteur de $\mathcal C$ sur 0 convient. \bookend

Soit alors $H$ un sous-espace strict de $F$, et $X$ vecteur de $F\setminus H$, on a $X\ne0$, (car $0\in H$), donc il existe une forme $\mu$ non nulle sur $X$, nulle sur $H$, d'où un $\mu$ non nul dans $H^\perp$ :
\[
H\subsetneq F\Longrightarrow H^\perp\ne\{0\}.
\]
Comme par ailleurs $F^\perp=\{\mu;\ \mu\in F^*,\ \mu=0\}$, c'est $F^\perp=\{0\}$, d'où finalement, pour $H$ sous-espace de $F$, on a
\[
H^\perp=\{0\}\Longleftrightarrow H=F.
\]
Il en résulte que
\[
u\text{ surjective}\Longleftrightarrow u(E)=F\Longleftrightarrow(u(E))^\perp=\{0\}
\Longleftrightarrow\operatorname{Ker}({}^tu)=\{0\}
\]
soit $u$ surjective $\Longleftrightarrow{}^tu$ injective. \bookend

\medskip\noindent\textsc{Théorème 6.105.} --- \emph{Soient $E$ et $F$ deux espaces vectoriels et $u\in\mathcal L(E,F)$. On a
\[
\operatorname{Im}({}^tu)=(\operatorname{Ker} u)^\perp.
\]}

Soit $\varphi\in\operatorname{Im}({}^tu)$, il existe $\psi$ dans $F^*$ tel que $\varphi={}^tu(\psi)=\psi\circ u$ donc, $\forall X\in\operatorname{Ker} u$,
\[
\varphi(X)=\psi(u(X))=\psi(0)=0:
\]
on a bien $\varphi\in(\operatorname{Ker} u)^\perp$.

Réciproquement soit $\varphi\in(\operatorname{Ker} u)^\perp$, on veut l'écrire sous la forme $\varphi={}^tu(\psi)=\psi\circ u$, avec $\psi$ dans $F^*$, c'est-à-dire factoriser $u$ à travers $\varphi$ or on sait (Théorème 6.54), que c'est possible si et seulement si $\operatorname{Ker} u\subset\operatorname{Ker}\varphi$, ce qui est le cas ici car $\forall X\in\operatorname{Ker} u$, $\varphi(X)=0$, donc $X\in\operatorname{Ker}\varphi$. \bookend

\medskip\noindent\textsc{Corollaire 6.106.} --- \emph{On a $u$ injective $\Longleftrightarrow{}^tu$ surjective.}

Car soit $E$ un espace vectoriel, on a $\{0\}^\perp=E^*$ car tout $\varphi$ de $E^*$ annule bien 0, et si $H$ est sous-espace $\ne\{0\}$ de $E$, avec $X\ne0$ dans $H$, il existe une forme coordonnée $\mu$ non nulle en $X$, (lemme 6.104), donc $\mu\notin H^\perp$ qui est donc différent de $E^*$. Il en résulte qu'avec $H$ sous-espace vectoriel de $E$ on a
\[
H^\perp=E^*\Longleftrightarrow H=\{0\},
\]
et donc ici, que
\[
u\text{ injective}\Longleftrightarrow\operatorname{Ker} u=\{0\}\Longleftrightarrow(\operatorname{Ker} u)^\perp=E^*
\Longleftrightarrow\operatorname{Im}({}^tu)=E^*
\]
soit $u$ injective $\Longleftrightarrow{}^tu$ surjective. \bookend

Il en résulte que l'on a :

\medskip\noindent\textsc{Théorème 6.107.} --- \emph{Soit $E$ un espace vectoriel sur un corps $K$, on a, pour $u$ application linéaire de $E$ dans $E$,
\[
u\in GL(E)\Longleftrightarrow{}^tu\in GL(E^*).
\]}

C'est une conséquence immédiate des corollaires 6.103 et 6.106. \bookend

Et les bases dans tout cela?

Soit $\mathcal B=(e_i)_{i\in I}$ une base de l'espace vectoriel $E$, comme $K$ est de dimension finie, 1, sur lui-même, il résulte du théorème 6.78, que les applications $(e_i^*)_{i\in I}$, linéaires de $E$ dans $K$ définies par les conditions :
\[
e_i^*(e_j)=0\quad\text{si }j\ne i\qquad\text{et}\qquad e_i^*(e_i)=1
\]
forment une famille libre de $E^*$, (car $\{1\}$ est une base de $K$ espace vectoriel sur lui-même).

Cette famille est une base de $E^*\Longleftrightarrow\dim E$ finie, (toujours d'après le Théorème 6.78) et dans ce cas, si $n=\dim E$, on a aussi $\dim E^*=n$.

\medskip\noindent\textbf{6.108.} Dans ce cas $\{e_1^*,\ldots,e_n^*\}$ est la base duale de $E^*$, mais il faut bien noter qu'en dimension infinie, la famille des $(e_i^*)_{i\in I}$ est libre mais non génératrice.

\medskip\noindent\textsc{Remarque 6.109.} --- Si
\[
X=\sum_{j\in I}x_j e_j
\]
avec les $x_j$ presque tous nuls, on a
\[
e_i^*(X)=\sum_{j\in I}x_j e_i^*(e_j)
\]
par linéarité, avec, pour $j\ne i$, $e_i^*(e_j)=0$ et $e_i^*(e_i)=1$, il reste $e_i^*(X)=x_i$ : $e_i^*$ est l'application qui à $X$ associe sa coordonnée suivant $e_i$, dans sa décomposition dans la base $\mathcal B$ d'où le nom de forme coordonnée.

Quelques résultats en dimension finie qui peuvent conduire à des justifications plus simples des théorèmes 6.96, 6.99, et 6.102 et 6.105, résultats non valables en dimension infinie.

\medskip\noindent\textsc{Théorème 6.110.} --- \emph{Soit $F$ sous-espace vectoriel de dimension $p$, de $E$ espace de dimension $n$. On a
\[
\dim F^\perp=n-p.
\]}

Car si $\{e_1,\ldots,e_p\}$ est une base de $F$, que l'on complète en
\[
\mathcal B=\{e_1,\ldots,e_p,e_{p+1},\ldots,e_n\}
\]
base de $E$, avec $\mathcal B^*$ base duale de $E^*$, on a
\[
\varphi=\sum_{i=1}^n\lambda_i e_i^*\in F^\perp=(\operatorname{Vect}(e_1,\ldots,e_p))^\perp
\]
si et seulement si pour tout $j\le p$, $\varphi(e_j)=0$ soit
\[
\sum_{i=1}^n\lambda_i e_i^*(e_j)=0.
\]
Il ne reste que $\lambda_j=0$, $\forall j=1,\ldots,p$, ce qui équivaut à
\[
\varphi=\sum_{j=p+1}^n\lambda_j e_j^*
\]
donc $F^\perp=\operatorname{Vect}(e_{p+1}^*,\ldots,e_n^*)$ est de dimension $n-p$. \bookend

Voyons, à titre d'exemple, comment la justification du Théorème 6.96, 2e partie est simplifiée, si $E$ est de dimension finie, $n$.

On veut prouver que $(F\cap G)^\perp=F^\perp+G^\perp$. On vérifie facilement que
\[
F^\perp+G^\perp\subset(F\cap G)^\perp,
\]
(car $(F\cap G)\subset F$ et $(F\cap G)\subset G$ impliquent $F^\perp\subset(F\cap G)^\perp$ et $G^\perp\subset(F\cap G)^\perp$ d'où $(F^\perp+G^\perp)\subset(F\cap G)^\perp$).

Puis
\[
\begin{aligned}
\dim(F^\perp+G^\perp)
&=\dim F^\perp+\dim G^\perp-\dim(F^\perp\cap G^\perp)\\
&=n-\dim F+n-\dim G-\dim(F+G)^\perp.
\end{aligned}
\]
La première partie du Théorème 6.96 donne $(F+G)^\perp=F^\perp\cap G^\perp$. On a donc
\[
\begin{aligned}
\dim(F^\perp+G^\perp)
&=2n-\dim F-\dim G-(n-\dim(F+G))\\
&=n-\dim F-\dim G+(\dim F+\dim G-\dim F\cap G)\\
&=n-\dim(F\cap G)=\dim(F\cap G)^\perp,
\end{aligned}
\]
(on utilise abondamment le corollaire 6.87). Mais la remarque 6.88 prouve l'égalité des deux sous-espaces, (inclusion et dimensions égales, finies). \bookend

\medskip\noindent\textsc{Théorème 6.111.} --- \emph{Soient $E$ et $F$ deux espaces vectoriels de dimensions finies $n$ et $p$ respectivement, et $u\in\mathcal L(E,F)$. Alors $u$ et ${}^tu$ ont le même rang.}

Car si $r=\operatorname{rang}u$, si $\{Y_1,\ldots,Y_r\}$ est une base de $u(E)$, si on choisit $X_1,\ldots,X_r$ dans $E$ tels que $u(X_i)=Y_i$, ($i=1,\ldots,r$), la famille $\{X_1,\ldots,X_r\}$ est libre, (si elle était liée, les images par $u$ le seraient).

On complète en $\{Y_1,\ldots,Y_r,Z_1,\ldots,Z_{p-r}\}$ base de $F$ cette famille libre des $Y_i$.

Soit
\[
\mathcal B^*=\{Y_1^*,\ldots,Y_r^*,Z_1^*,\ldots,Z_{p-r}^*\}
\]
la base duale de $F^*$, le rang de ${}^tu$ est la dimension de l'espace vectoriel engendré par les ${}^tu(Y_i^*)$ et les ${}^tu(Z_k^*)$. Or ${}^tu(Z_k^*)=Z_k^*\circ u=0$ car pour tout $X$ de $E$, $u(X)\in\operatorname{Vect}(Y_1,\ldots,Y_r)$ donc a sa coordonnée suivant $Z_k$ nulle, $Z_k^*(u(X))=$ coordonnée de $u(X)$ suivant $Z_k$, d'après 6.109. Il reste
\[
\operatorname{rang}({}^tu)=\dim\operatorname{Vect}({}^tu(Y_1^*),\ldots,{}^tu(Y_r^*)),
\]
mais ces formes, égales à $Y_1^*\circ u,Y_2^*\circ u,\ldots,Y_r^*\circ u$ sont libres car si on a des scalaires $\lambda_1,\ldots,\lambda_r$ tels que
\[
\sum_{i=1}^r\lambda_iY_i^*\circ u=0,
\]
en prenant l'image de $X_j$ on trouve 0, or $Y_i^*(u(X_j))=Y_i^*(Y_j)=0$ sauf si $i=j$, et $Y_j^*(Y_j)=1$. Il reste $\lambda_j=0$, ceci pour $j=1,\ldots,r$.

Finalement ${}^tu(F^*)$ est bien de dimension $r$, d'où $\operatorname{rang}({}^tu)=r$. \bookend

En fait ce théorème reste valable, quelles que soient les dimensions de $E$ et $F$, dès que $u$ est de rang fini. On a le :

\medskip\noindent\textsc{Théorème 6.112.} --- \emph{Soient $E$ et $F$ deux espaces vectoriels, $u\in\mathcal L(E,F)$. Si $r=\operatorname{rang}u=\dim u(E)$ est fini alors ${}^tu$ est aussi de rang fini, $r$.}

On suppose $u$ de rang fini $r$, soit $\{Y_1,\ldots,Y_r\}$ une base de $u(E)$ et des vecteurs $(X_i)_{1\le i\le r}$ de $E$ tels que $u(X_i)=Y_i$. En considérant la justification du théorème 6.111, on s'aperçoit qu'elle reste valable si $F$ est de dimension finie : la dimension de $E$ n'intervenant pas en fait.

Donc on suppose ici $F$ de dimension infinie. On complète la famille libre $\{Y_1,\ldots,Y_r\}$ en
\[
\mathcal B=\{Y_1,\ldots,Y_r\}\cup(Z_j)_{j\in J}
\]
base de $F$ (théorème de la base incomplète), d'où, dans $F^*$, la famille libre notée $\mathcal L$ formée des formes coordonnées $(Y_i^*)_{1\le i\le r}$ et $(Z_j^*)_{j\in J}$ qui envoient chacune toute la base sur 0 sauf un vecteur d'image 1.

Cette famille libre, n'est pas une base de $F^*$, (6.108) car $F$ est de dimension infinie mais il existe $\mathcal C$ famille libre de $F^*$ telle que $\mathcal L\cup\mathcal C$ soit base de $F^*$, là encore théorème de la base incomplète.

L'image ${}^tu(F^*)$ est alors engendrée par les images des vecteurs de cette base $\mathcal L\cup\mathcal C=(Y_i^*)_{1\le i\le r}\cup(Z_j^*)_{j\in J}\cup\mathcal C$ par ${}^tu$.

On aura ${}^tu(Z_j^*)=0$ car c'est $Z_j^*\circ u$ et $\forall X\in E$, $u(X)\in u(E)=\operatorname{Vect}(Y_1,\ldots,Y_r)$, donc a des coordonnées nulles suivant les $Z_j$.

Mais \emph{a priori}, si $\varphi\in\mathcal C$, il n'y a aucune raison pour que ${}^tu(\varphi)=0$. Aussi la démonstration consiste à remplacer $\mathcal C$ par une famille libre $\mathcal C'$ formée de $\varphi'$ tels que ${}^tu(\varphi')=0$, soit $\varphi'\circ u=0$, ou $\varphi'$ nulle sur $u(E)$.

Pour cela, si $\varphi\in\mathcal C$, on définit $\varphi'$ associée par :
\[
\varphi'=\varphi-\sum_{i=1}^r\varphi(Y_i)Y_i^*.
\]

$\varphi'$ est bien un élément de $F^*$, et sur chaque vecteur $Y_k$, ($1\le k\le r$) de la base considérée de $u(E)$ on a
\[
\varphi'(Y_k)=\varphi(Y_k)-\sum_{i=1}^r\varphi(Y_i)Y_i^*(Y_k),
\]
avec $Y_i^*(Y_k)=0$ sauf si $i=k$, et alors $Y_k^*(Y_k)=1$, il reste donc $\varphi'(Y_k)=0$ : on a bien $\varphi'\circ u=0$.

Si on note $\mathcal C'=\{\varphi'\text{ ainsi associés aux }\varphi\text{ de }\mathcal C\}$ on a $\mathcal L\cup\mathcal C'$ base de $F^*$ car tout élément de $F^*$ est combinaison linéaire d'un nombre fini d'éléments de $\mathcal L\cup\mathcal C$. Si on remplace les $\varphi$ de $\mathcal C$ intervenant par les $\varphi'$ associés on garde une combinaison linéaire des mêmes éléments de $\mathcal L$, des $\varphi'$ en nombre fini, et en plus de $Y_1^*,\ldots,Y_r^*$ qui sont dans $\mathcal L$; finalement tout élément de $F^*$ est combinaison linéaire d'un nombre fini d'éléments de $\mathcal L\cup\mathcal C'$ qui est donc partie génératrice de $F^*$. C'est une famille libre, (donc une base) car supposons que l'on ait une combinaison linéaire finie, nulle des éléments de $\mathcal L\cup\mathcal C'$.

En notant $\mathcal C=(\varphi_l)_{l\in L}$, et en indexant $\mathcal C'$ par le même ensemble, et quitte à rajouter à la combinaison linéaire non nulle, 0 fois chaque $Y_i^*$ n'y figurant pas, on peut écrire cette combinaison sous la forme
\[
\sum_{i=1}^r y_iY_i^*+\sum_{a=1}^p z_{j_a}Z_{j_a}^*+\sum_{b=1}^q t_{l_b}\varphi'_{l_b}=0,
\]
car pour prendre un nombre fini, $p$, des $Z_j^*$, on indexe les indices $j$ par un indice, ici $a$, variant de 1 à $p$. On procède de même pour les $\varphi'_l$.

On remplace chaque $\varphi'_{l_b}$ par
\[
\varphi_{l_b}-\sum_{i=1}^r\varphi_{l_b}(Y_i)Y_i^*,
\]
il reste
\[
\sum_{i=1}^r\left(y_i-\sum_{b=1}^q t_{l_b}\varphi_{l_b}(Y_i)\right)Y_i^*
+\sum_{a=1}^p z_{j_a}Z_{j_a}^*
+\sum_{b=1}^q t_{l_b}\varphi_{l_b}=0,
\]
(c'est affreux, mais il faut savoir s'accrocher, comme le fou à son pinceau).

En tout cas, $y_i-\sum_{b=1}^q t_{l_b}\varphi_{l_b}(Y_i)$ est un scalaire, et là, c'est une combinaison linéaire des vecteurs de la base $\mathcal L\cup\mathcal C$ de $F^*$ donc les $t_{l_b}$ sont nuls, les $z_{j_a}$ aussi, il reste $\sum_{i=1}^r y_iY_i^*=0$ donc les $y_i$ aussi sont nuls : on a finalement une base $\mathcal L\cup\mathcal C'$ de $F^*$.

En revenant à la justification du théorème, l'image de ${}^tu$ est engendrée par
\begin{itemize}[leftmargin=2em]
\item les ${}^tu(Y_i^*)_{1\le i\le r}$,
\item les ${}^tu(Z_j^*)$ qui sont nuls,
\item et les ${}^tu(\varphi')$, $\forall\varphi'\in\mathcal C'$, mais on a construit les $\varphi'$ pour avoir 0.
\end{itemize}
Il reste
\[
{}^tu(F^*)=\operatorname{Vect}({}^tu(Y_i^*)_{1\le i\le r}),
\]
et on termine comme pour le théorème 6.111 cette famille est libre car si on a des scalaires $\lambda_1,\ldots,\lambda_r$ tels que
\[
\sum_{i=1}^r\lambda_i{}^tu(Y_i^*)=0,
\]
en prenant l'image de $X_k$ choisi tel que $u(X_k)=Y_k$, on a
\[
0=\left(\sum_{i=1}^r\lambda_i{}^tu(Y_i^*)\right)(X_k)
=\sum_{i=1}^r\lambda_iY_i^*(u(X_k))=\lambda_k,
\]
donc $\lambda_k$ nul.

Il en résulte que les $({}^tu(Y_i^*))_{1\le i\le r}$ forment une base de ${}^tu(F^*)$ donc ${}^tu$ est bien de rang $r$. Ouf! \bookend

Il reste une question à aborder.

Si $E$ est de dimension finie, $n$, son dual $E^*$ est aussi de dimension $n$. Si $E$ est de dimension infinie, et si $\mathcal B$ est une base de $E$, les formes coordonnées forment une famille libre de $E^*$, pas une base mais ... quelle est la dimension du dual $E^*$ ?

Une base de $E^*$ pourrait-elle être équipotente à une de $E$? La réponse à cette question passionnante est connue : si $E$ est de dimension infinie, son dual a une dimension strictement plus grande, ou encore une base de $E^*$ a un cardinal strictement supérieur \BellaCorrectionNote{au cardinal}{ou cardinal}{Correction de la préposition dans le texte imprimé.} d'une base de $E$.

C'est ce qui va être justifié par le théorème suivant :

\medskip\noindent\textsc{Théorème 6.113. (d'Erdös Kaplansky)} --- \emph{Soit $E$ un espace vectoriel de dimension infinie sur un corps $K$. La dimension de son dual $E^*$ est $\operatorname{card}(E^*)$.}

Oui, cela surprend! Cela ne signifie pas que $E^*$ est une base de $E^*$, ce qui serait ridicule, mais qu'une base de $E^*$ est équipotente à $E^*$, ce qui ne l'empêche pas d'avoir «moins d'éléments» au sens intuitif du «moins» qui dit qu'il y a «moins» d'entiers que de rationnels bien qu'il y en ait « autant » puisque ces deux ensembles sont équipotents.

On comprend que la justification ne va pas être facile, et s'il y a encore des lecteurs, qu'ils s'accrochent bien pour se plonger dans les délices des infinis. A moins que, raisonnablement, ils passent à autre chose.

Comme une base $\mathcal B$ de $E^*$ est contenue dans $E^*$, on a déjà
\[
\dim(E^*)=\operatorname{card}\mathcal B\le\operatorname{card}(E^*),
\]
il suffit donc de justifier l'autre inégalité. On peut considérer qu'une moitié du travail est faite, la plus petite.

Soit $\mathcal C=(e_i)_{i\in I}$ une base de $E$. Comme une forme linéaire $\varphi$ est déterminée par les valeurs prises sur les $e_i$, on a :

\medskip\noindent\textsc{Lemme 6.114.} --- \emph{L'espace vectoriel $E^*$ est isomorphe à $K^I$ espace vectoriel des applications de $I$ dans $K$.}

On définit l'application $\theta$ de $E^*$ dans $K^I$ par : $\theta(\varphi)$ est l'application $i\longmapsto\varphi(e_i)$ de $I$ dans $K$.

$\theta$ linéaire : si $\varphi$ et $\psi$ sont dans $E^*$, $\lambda$ et $\mu$ dans $K$, $\theta(\lambda\varphi+\mu\psi)$ est l'application de $I$ dans $K$ qui à $i$ associe $(\lambda\varphi+\mu\psi)(e_i)$ soit
\[
\lambda\varphi(e_i)+\mu\psi(e_i)=\lambda\theta(\varphi)(i)+\mu\theta(\psi)(i)
=[\lambda\theta(\varphi)+\mu\theta(\psi)](i).
\]
Ceci étant vrai pour tout $i$ de $I$ c'est que $\theta(\lambda\varphi+\mu\psi)$ est bien $\lambda\theta(\varphi)+\mu\theta(\psi)$.

$\theta$ injective : si $\theta(\varphi)=0$, la forme $\varphi$ valant 0 sur la base $\mathcal C$ de $E$ est nulle.

$\theta$ surjective : car la donnée d'une application $a$ de $I$ dans $K$ détermine une forme $\varphi$ valant $a(i)$ sur $e_i$ (Théorème 6.77). \bookend

Par ailleurs, $(e_i)_{i\in I}$ étant base de $E$, les vecteurs de $E$ sont identifiés aux familles $(a_i)_{i\in I}$ de scalaires presque tous nuls, on va noter $K^{[I]}$ l'espace vectoriel des applications de $I$ dans $K$ à valeurs presque toutes nulles, c'est-à-dire telles que $\operatorname{card}\{i,a_i\ne0\}$ est fini. Le fait que $K^{[I]}$ soit un espace vectoriel, sous-espace de $K^I$ en fait, est facile à vérifier.

\medskip\noindent\textsc{Lemme 6.115.} --- \emph{Soit $I$ un ensemble, $a^{(1)},a^{(2)},\ldots,a^{(k)}$ des applications distinctes en nombre fini dans $K$. Alors il existe une partie $J$ de cardinal fini, $J\subset I$, telle que les éléments $(a_i^{(p)})_{i\in J}$ de $K^{\operatorname{card}J}$ soient distincts, (pour $p=1,2,\ldots,k$).}

Digérons d'abord les notations. Une application $a$ de $I$ dans $K$ est ici notée comme une famille $(a_i)_{i\in I}$, $a_i$ étant la valeur prise par la fonction $a$ en $i$. C'est la notation utilisée pour les suites que l'on généralise ici.

Comme on considère plusieurs applications, on les indexe par un indice $p$ variant de 1 à $k$, et placé en haut à droite, entre parenthèses pour ne pas confondre avec des puissances.

Enfin si $\operatorname{card}(J)=n$ est un entier, $K^n$ est un espace vectoriel sur $K$, celui des $n$-uplets, (exemple 6.3).

Soient alors $p$ et $q$ distincts dans $\{1,2,\ldots,k\}$, les fonctions $a^{(p)}$ et $a^{(q)}$ étant distinctes, il existe un élément $j_{p,q}$ de $I$ où elles prennent des valeurs différentes, donc
\[
a_{j_{p,q}}^{(p)}\ne a_{j_{p,q}}^{(q)}.
\]
Si
\[
J=\bigcup_{\{p,q\}\text{ paire de }\{1,2,\ldots,k\}}\{j_{p,q}\}
\]
est un ensemble formé d'indices ainsi associés aux $\frac{k(k-1)}2$ sous-ensembles à 2 éléments de $\{1,\ldots,k\}$ et si $n=\operatorname{card}J$, on a $n\le\frac{k(k-1)}2$, et les $n$-uplets $(a_i^{(p)})_{i\in J}$ pour $p=1,2,\ldots,k$ sont bien distincts. \bookend

\medskip\noindent\textsc{Lemme 6.116.} --- \emph{Soit $n$ un entier $\ge1$, $b^{(1)},b^{(2)},\ldots,b^{(k)}$ des $n$-uplets distincts. Il existe un polynôme $P$ à $n$ indéterminées $X_1,\ldots,X_n$ sur $K$, tel que $P(b^{(1)})\ne0$ et $P(b^{(s)})=0$ si $2\le s\le k$.}

Rappelons que si
\[
P(X_1,\ldots,X_n)=\sum_{(j_1,j_2,\ldots,j_n)\in\mathbb N^n}u_{j_1\cdots j_n}X_1^{j_1}X_2^{j_2}\cdots X_n^{j_n}
\]
avec les scalaires $u_{j_1\cdots j_n}$ de $K$, indexés par $\mathbb N^n$, presque tous nuls, et si le $n$-uplet $b^{(s)}$ de $K^n$ vaut $b^{(s)}=(b_1^{(s)},b_2^{(s)},\ldots,b_n^{(s)})$ on a
\[
P(b^{(s)})=\sum_{j_1,\ldots,j_n}u_{j_1\cdots j_n}(b_1^{(s)})^{j_1}(b_2^{(s)})^{j_2}\cdots(b_n^{(s)})^{j_n}.
\]

On justifie le lemme par récurrence sur $k$.

Si $k=2$, on a $b^{(1)}\ne b^{(2)}$ donc il existe un indice $j$ dans $\{1,\ldots,n\}$ tel que $b_j^{(1)}\ne b_j^{(2)}$, le polynôme
\[
P(X_1,\ldots,X_n)=X_j-b_j^{(2)}
\]
s'annule sur le $n$-uplet $b^{(2)}$ et pas sur $b^{(1)}$, car $P(b^{(1)})=b_j^{(1)}-b_j^{(2)}\ne0$.

Si on suppose le résultat vrai à l'ordre $k$, soient $b^{(1)},\ldots,b^{(k)},b^{(k+1)}$ des points distincts de $K^n$, on sait (hypothèse de récurrence) qu'il existe un polynôme $Q$ en $X_1,\ldots,X_n$, tel que $Q(b^{(1)})\ne0$ et $Q(b^{(s)})=0$ pour $s=2,3,\ldots,k$. Puis, (cas de $k=2$), il existe un polynôme en $X_1,\ldots,X_n$ également, $R$, tel que $R(b^{(1)})\ne0$ et $R(b^{(k+1)})=0$. Alors
\[
P(X_1,\ldots,X_n)=Q(X_1,\ldots,X_n)R(X_1,\ldots,X_n)
\]
est tel que $P(b^{(1)})=Q(b^{(1)})R(b^{(1)})\ne0$, et pour $s=2,\ldots,k,k+1$, $P(b^{(s)})=0$ car soit $Q(b^{(s)})=0$ soit $R(b^{(s)})=0$. \bookend

Avant de continuer, un petit peu de gymnastique mentale.

Soit $I$ un ensemble, $S=\{X_i,i\in I\}$ un ensemble d'indéterminées indexées par $I$, donc en bijection avec $I$.

\medskip\noindent\textbf{6.117.} On introduit $A=K[S]$ l'ensemble des polynômes ayant $S$ comme ensemble d'indéterminées.

On a donc
\[
\begin{aligned}
P\in A\quad\Longleftrightarrow\quad &\exists I'\text{ partie finie de }I,\\
& P\text{ est un polynôme par rapport aux indéterminées }X_i,\ i\in I'.
\end{aligned}
\]

Si $Q\in A$ est polynôme par rapport aux $X_i$ avec $i\in I''$, et $\operatorname{card}I''$ fini, avec $I'''=I'\cup I''$, de cardinal fini, on a aussi $P$ et $Q$ polynômes par rapport aux indéterminées $X_i$, en nombre fini, associées aux $i$ de $I'''$, donc on peut définir $P+Q$ et $PQ$ avec les règles usuelles, ainsi que $\lambda\cdot P$ si $\lambda$ est un scalaire, si bien qu'avec un peu de réflexion on parvient à concevoir que $A$ est une algèbre.

Revenons à nos moutons, avec $E$ vectoriel de dimension infinie, $\mathcal C=(e_i)_{i\in I}$ une base de $E$, $S=(X_i)_{i\in I}$ une famille d'indéterminées indexées par $I$ et $A$ l'algèbre des polynômes aux indéterminées $(X_i)_{i\in I}$.

\medskip\noindent\textsc{Lemme 6.118.} --- \emph{Soit $a$ dans $K^I$, une application de $I$ dans $K$, (on notera $a_i$ la valeur prise par $a$ en $i$). Soit $\delta_a$ l'application de $A=K[S]$ dans $K$ définie par $\delta_a(P)=P(a)$. La famille $(\delta_a)_{a\in K^I}$ est libre dans $A^*$, dual de $A$.}

Un point pour comprendre, si tant est que cela soit encore possible! Si $P$ est un polynôme de $K[S]$, par rapport aux variables $X_{i_1},\ldots,X_{i_r}$ en nombre fini, (les $i_k\in I$), $P(a)$ désigne la valeur prise par $P$ quand on remplace les indéterminées $X_{i_1},\ldots,X_{i_r}$ par $a_{i_1},\ldots,a_{i_r}$.

Retour au lemme. Soient $a^{(1)},a^{(2)},\ldots,a^{(k)}$ des éléments distincts en nombre fini de $K^I$, on va prouver que les applications $\delta_{a^{(p)}}$, $p=1,\ldots,k$, sont des éléments indépendants de $A^*$.

D'abord ce sont des éléments de $A^*$ dual de $A$ car l'application qui pour $a$ fixé dans $K^I$, associe $P(a)$ au polynôme $P$ variant dans $A$ est bien linéaire en $P$, car
\[
(\lambda P+\mu Q)(a)=\lambda P(a)+\mu Q(a)
\]
par définition de la structure vectorielle sur $A$. De plus c'est à valeurs dans le corps de base.

Si on a des scalaires $\lambda_1,\ldots,\lambda_k$ tels que
\[
\sum_{p=1}^k\lambda_p\delta_{a^{(p)}}=0,
\]
il faut prouver qu'ils sont tous nuls.

Procédons par l'absurde : on les suppose non tous nuls, par exemple $\lambda_1\ne0$. Le lemme 6.115 nous dit qu'il existe une partie $J$ de $I$, de cardinal fini, $n$, et telle qu'en posant
\[
b^{(p)}=(a_i^{(p)})_{i\in J},
\]
on obtient $b^{(1)},\ldots,b^{(k)}$ qui sont $k$ éléments distincts de $K^n$.

D'après le lemme 6.116, il existe alors un polynôme par rapport aux $(X_i)_{i\in J}$ donc $P\in A$ vu la définition de $A$, tel que $P(b^{(1)})\ne0$ et $P(b^{(s)})=0$ pour $s=2,3,\ldots,k$. Mais il est clair que l'on a
\[
P(b^{(s)})=P(a^{(s)})=\delta_{a^{(s)}}(P),
\]
d'où $\delta_{a^{(1)}}(P)\ne0$ et $\delta_{a^{(s)}}(P)=0$, pour $s=2,\ldots,k$.

(Dans $P(b^{(s)})$, $P$ est considéré comme un polynôme en $n$ variables, et on le calcule en $b^{(s)}$ qui est un $n$-uplet, alors que dans $P(a)$ on considère $P$ en tant qu'élément de $A$.)

On a donc
\[
\left(\sum_{p=1}^k\lambda_p\delta_{a^{(p)}}\right)(P)=\lambda_1\delta_{a^{(1)}}(P)\ne0
\]
avec $\lambda_1\ne0$ : c'est absurde. \bookend

Il nous reste un dernier lemme à justifier, avant de passer au théorème d'Erdös Kaplansky.

\medskip\noindent\textsc{Lemme 6.119.} --- \emph{Soit $E$ un espace vectoriel de dimension infinie, $\mathcal C=(e_i)_{i\in I}$ une base de $E$ et $S=(X_i)_{i\in I}$ une famille d'indéterminées indexée par $I$. Alors $E$ est isomorphe à l'espace vectoriel $A=K[S]$ des polynômes par rapport aux indéterminées $(X_i)_{i\in I}$.}

En effet
\[
\begin{aligned}
\mathcal M&=\{\text{monômes unitaires par rapport aux }X_i\}\\
&=\{1\}\cup\Bigl\{X_{i_1}^{p_1}\cdots X_{i_k}^{p_k};\ \{i_1,\ldots,i_k\}\text{ partie finie de }I,\\
&\hspace{8.5em}p_1,\ldots,p_k\in\mathbb N^*\Bigr\}.
\end{aligned}
\]
est une base de $A$ car un polynôme est combinaison linéaire de monômes, et une combinaison linéaire de monômes est identiquement nulle si et seulement si chaque coefficient est nul.

Il existe une injection de $I$ dans $\mathcal M$ : tout simplement $i\longmapsto X_i$ donc $\operatorname{card}\mathcal M\ge\operatorname{card}I$.

En ne considérant que $\mathcal M'$, l'ensemble des monômes autres que 1, de tels monômes sont caractérisés par la donnée de la partie $\{i_1,\ldots,i_k\}$ des indices des variables, partie finie non vide de $I$, et des valeurs $p_1,\ldots,p_k$ dans $\mathbb N^*$. De plus, $\mathcal M$ est de cardinal infini, comme $I$, donc $\mathcal M'$ aussi puisque $\mathcal M'\cup\{1\}=\mathcal M$, et $\operatorname{card}(\mathcal M')=\operatorname{card}(\mathcal M)$.

Soit $J=\{i_1,\ldots,i_k\}$ une partie finie de $k$ éléments de $I$. Les applications de $J$ dans $\mathbb N^*$ sont les éléments de $\mathbb N^*\times\mathbb N^*\times\cdots\times\mathbb N^*$, ($k$ fois), le cardinal de cet ensemble est $(\operatorname{card}\mathbb N)^k=\operatorname{card}\mathbb N$, (proposition 3.52), il existe donc une bijection $\theta_J$ de l'ensemble des applications de $J$ dans $\mathbb N^*$ sur l'ensemble $\mathbb N$.

On considère alors l'ensemble $\mathcal F$ des sous-familles finies $(i_1,\ldots,i_k)$ de $I$. L'application $\theta$ qui au monôme $X_{i_1}^{p_1}\cdots X_{i_k}^{p_k}$ associe le couple
\[
(J,\theta_J(p_1,\ldots,p_k))\quad\text{de }\mathcal F\times\mathbb N
\]
avec $J=(i_1,\ldots,i_k)$ et $\theta_J(p_1,\ldots,p_k)$ le nombre entier associé par $\theta_J$ à l'application $i_1\mapsto p_1;\ldots;i_k\mapsto p_k$ est une bijection de $\mathcal M'$ sur $\mathcal F\times\mathbb N$, vu la caractérisation des monômes.

Or le cardinal de l'ensemble $\mathcal F$ des parties finies de $I$, de cardinal infini, est encore $\operatorname{card}(I)$. On a vu dans la justification du théorème 6.74, que si $A_k$ est l'ensemble des parties de $k$ éléments de $I$, on a $\operatorname{card}A_k=\operatorname{card}I$ d'où une bijection $\varphi_k$ de $A_k$ sur $I$ d'où une, $\varphi$, de $\mathcal F$ sur $\mathbb N^*\times I$, celle qui à la partie finie non vide $J=\{i_1,\ldots,i_k\}$ va associer le couple $(k,\varphi_k(J))$. Donc $\mathcal F$ et $\mathbb N^*\times I$ sont équipotents,
\[
\operatorname{card}(\mathcal F)=\operatorname{card}(\mathbb N^*)\times\operatorname{card}(I)=\operatorname{card}(I)
\]
d'après la proposition 3.52, $\operatorname{card}(\mathbb N^*)$ étant le plus petit cardinal infini.

On a finalement
\[
\operatorname{card}(\mathcal M')=\operatorname{card}(\mathcal F)\times\operatorname{card}(\mathbb N)=\operatorname{card}(I),
\]
et $A$ et $E$ ont bien des bases équipotentes : ils sont isomorphes. \bookend

Si $u$ est un isomorphisme de $E$ sur $A$, ${}^tu$ en est un de $A^*$ sur $E^*$ (théorème 6.107) qui sont donc isomorphes : $E^*$ et $A^*$ ont même dimension, infinie.

Mais la famille $(\delta_a)_{a\in K^I}$ est libre dans $A^*$, (lemme 6.118), donc
\[
\dim E^*=\dim A^*\ge\operatorname{card}((\delta_a)_{a\in K^I})=\operatorname{card}(K^I)
\]
et $K^I$ étant isomorphe à $E^*$, (lemme 6.114), ces 2 ensembles sont équipotents d'où l'inégalité $\dim(E^*)\ge\operatorname{card}(E^*)$ voulue, ce qui achève la justification du théorème d'Erdös Kaplansky. On a bien, pour $E$ de dimension infinie, $E^*$ de dimension $\operatorname{card}(E^*)$ aussi bizarre que cela paraisse. \bookend

\medskip\noindent\textsc{Corollaire 6.120.} --- \emph{Pour $E$, vectoriel de dimension infinie, $E$ et $E^*$ ne sont jamais isomorphes.}

En effet, $E^*$ a ses bases de cardinal, celui de $E^*$, donc encore de cardinal celui de $K^I$. Or $K$ admet au moins 2 éléments, 0 et 1 les éléments neutres des 2 lois, donc $\{0,1\}^I$, ensemble des applications de $I$ dans $\{0,1\}$ s'injecte canoniquement dans $K^I$, ensemble des applications de $I$ dans $K$, ce qui donne
\[
\operatorname{card}(K^I)\ge\operatorname{card}(\{0,1\}^I).
\]
Or on aurait pu établir, dans le chapitre sur les cardinaux, le théorème suivant :

\medskip\noindent\textsc{Théorème 6.121.} --- \emph{Pour tout ensemble $I$ on a
\[
\operatorname{card}\{0,1\}^I=2^{\operatorname{card}I}>\operatorname{card}I,
\]
que nous justifierons ensuite.}

Si on l'applique ici, on obtient finalement
\[
\dim(E^*)=\operatorname{card}(K^I)\ge2^{\operatorname{card}(I)}>\operatorname{card}(I)=\dim(E):
\]
$E$ et $E^*$ ne sont pas isomorphes puisque leurs bases ne sont pas équipotentes. \bookend

Retour au théorème 6.121. Par convention, $2^0=1$, c'est bien $>0$.

Si $I\ne\varnothing$, il existe une injection de $I$ dans $\{0,1\}^I$, celle qui à $i$ de $I$ associe l'application $\chi_i$ définie par $\chi_i(i)=1$ et, $\forall j\ne i$, $j\in I$, $\chi_i(j)=0$. Donc
\[
\operatorname{card}(I)\le\operatorname{card}(\{0,1\}^I).
\]

Il n'y a pas de bijection $\varphi$ de $I$ sur $\{0,1\}^I$ car si $\varphi$ en est une, on définit $f$ de $I$ dans $\{0,1\}$ comme suit.

Pour $i\in I$, soit $\varphi(i)$ l'application de $I$ dans $\{0,1\}$ associée à $i$.

Si $\varphi(i)$ envoie $i$ sur 1 on pose $f(i)=0$,

si $\varphi(i)$ envoie $i$ sur 0 on pose $f(i)=1$.

L'application $f$ de $I$ dans $\{0,1\}$ ainsi définie n'est jamais image par $\varphi$ d'un élément de $I$ car si on suppose l'existence de $i_0$ dans $I$ tel que $\varphi(i_0)=f$, la valeur prise par $f$ en $i_0$ est forcément différente de celle prise par l'application $\varphi(i_0)$ en $i_0$.

Donc l'inégalité $\operatorname{card}(I)\le\operatorname{card}(\{0,1\}^I)$ est stricte ce qui justifie le théorème. \bookend
